\section{Urgensi Kriptografi dalam Teknologi Informasi}

\subsection{Mengapa Komunikasi Publik Rentan?}

Hampir semua komunikasi modern---surat elektronik, SMS, browser web, perbankan daring---melewati kanal publik (internet, jaringan seluler, infrastruktur bersama). Tanpa proteksi kriptografi, pesan dapat disadap, diubah, atau dipalsukan. Karena itu kriptografi menjadi inti banyak teknologi keamanan siber: memahami perannya penting untuk memahami cara mengamankan ruang siber \cite{munir_slides_itb,stallings2017}.

\subsection{Kasus Ilustratif: Kebocoran dan Serangan}

Beberapa kasus berikut bersifat ilustratif (bukan laporan forensik lengkap). Tujuannya menunjukkan bahwa kegagalan melindungi data berdampak nyata pada privasi, kepercayaan publik, dan ekonomi digital \cite{munir_slides_itb}:
\begin{itemize}
  \item \textbf{Kebocoran data pribadi berulang}: data peserta asuransi kesehatan, pengguna e-commerce, atau nasabah lembaga keuangan yang bocor ke publik menekankan pentingnya enkripsi, kontrol akses, dan tata kelola data.
  \item \textbf{Penyadapan atau akses tidak sah terhadap fasilitas komunikasi}: menggarisbawahi kebutuhan confidentiality pada kanal diplomatik maupun korporat.
  \item \textbf{Peretasan situs lembaga atau ransomware}: dokumen rahasia atau data nasabah berisiko diekspos atau dienkripsi paksa oleh penyerang---menunjukkan bahwa keamanan tidak hanya soal firewall, tetapi juga proteksi kriptografi dan manajemen kunci.
\end{itemize}

\begin{konsep}
Pelajaran utama: tanpa kerahasiaan, integritas, dan otentikasi yang memadai, data di kanal publik maupun di penyimpanan mudah menjadi sasaran. Kriptografi bukan satu-satunya solusi keamanan, tetapi merupakan fondasi teknis yang sulit digantikan.
\end{konsep}

\subsection{Enkripsi At Rest dan In Transit}

Dua aplikasi utama enkripsi dalam praktik TI \cite{munir_slides_itb}:
\begin{itemize}
  \item \textbf{Encryption at rest}: melindungi data yang disimpan (disk, basis data, berkas dokumen) agar tetap rahasia jika media fisik atau salinan cadangan dicuri.
  \item \textbf{Encryption in transit} (\textit{on motion}): melindungi pesan
        yang dikirim melalui jaringan.
        Contoh familiar adalah enkripsi ujung-ke-ujung pada aplikasi pesan
        atau HTTPS/TLS pada browser.
\end{itemize}

Gambar~\ref{fig:at-rest-in-transit} memberi contoh keduanya dari aplikasi sehari-hari.

\begin{figure}[htbp]
\centering
\begin{minipage}[b]{0.42\textwidth}
  \centering
  \includegraphics[width=\linewidth]{bab-02/encrypt-document}\\[2pt]
  {\small (a) At rest: enkripsi berkas dokumen dengan kata sandi}
\end{minipage}\hfill
\begin{minipage}[b]{0.48\textwidth}
  \centering
  \includegraphics[width=\linewidth]{bab-02/whatsapp-e2e}\\[2pt]
  {\small (b) In transit: enkripsi ujung-ke-ujung pada aplikasi pesan}
\end{minipage}
\caption{Contoh encryption at rest dan encryption in transit \cite{munir_slides_itb}.}
\label{fig:at-rest-in-transit}
\end{figure}

\begin{catatan}
Sistem modern sering menggabungkan keduanya: data dienkripsi saat disimpan \textit{dan} saat dikirim. TLS (Bab~VIII) adalah contoh proteksi in transit; mode penyimpanan seperti XTS (Bab~V) berkaitan dengan proteksi at rest.
\end{catatan}

\subsection{Lembaga Terkait Kriptografi di Indonesia}

Di Indonesia, pengembangan kapasitas sandi dan keamanan siber melibatkan lembaga negara dan pendidikan antara lain \cite{munir_slides_itb,bssn_official}:
\begin{enumerate}
  \item \textbf{Badan Siber dan Sandi Negara (BSSN)} --- \url{https://bssn.go.id} --- lembaga yang menggabungkan fungsi sandi negara dan aspek aplikasi informatika terkait keamanan siber.
  \item \textbf{Politeknik Siber dan Sandi Negara (PSSN)} --- \url{https://poltekssn.ac.id} --- pendidikan tinggi kedinasan di bidang siber dan sandi.
\end{enumerate}

Museum Sandi di Yogyakarta (\url{https://museum.lemsaneg.go.id/}, Jl.~Faridan Muridan Noto No.~21, Kota Baru) menyimpan koleksi alat sandi yang pernah digunakan di Indonesia dan menjadi referensi sejarah praktis bagi mahasiswa yang ingin melihat artefak kriptografi klasik secara langsung \cite{munir_slides_itb}. Salah satu koleksinya adalah mesin sandi SN-101 (Gambar~\ref{fig:museum-sandi}).

\begin{figure}[htbp]
\centering
\begin{minipage}[b]{0.48\textwidth}
  \centering
  \includegraphics[width=\linewidth]{bab-02/museum-sandi}\\[2pt]
  {\small (a) Gedung Museum Sandi, Yogyakarta}
\end{minipage}\hfill
\begin{minipage}[b]{0.44\textwidth}
  \centering
  \includegraphics[width=\linewidth]{bab-02/mesin-sandi-sn101}\\[2pt]
  {\small (b) Mesin sandi SN-101}
\end{minipage}
\caption{Museum Sandi dan salah satu koleksi alat sandinya \cite{munir_slides_itb}.}
\label{fig:museum-sandi}
\end{figure}

Bagian berikutnya merinci \textbf{layanan keamanan informasi} yang menjawab empat masalah Alice--Bob dan ancaman yang telah dibahas.
